\subsubsection{Comparação entre BJT e MOSFET}

Embora ambos os dispositivos possam ser utilizados como amplificadores e chaves eletrônicas, existem diferenças importantes em seus mecanismos de funcionamento.

\begin{table}[H]
    \centering
    \caption{Comparação entre transistores BJT e MOSFET}
    \label{tab:bjt-mosfet}
    \begin{tabular}{|l|c|c|}
        \hline
        \textbf{Característica} & \textbf{BJT} & \textbf{MOSFET} \\
        \hline
        Controle & Corrente & Tensão \\
        \hline
        Terminais & Emissor, base, coletor & Gate, source, drain \\
        \hline
        Impedância de entrada & Menor & Muito alta \\
        \hline
        Corrente de controle & Necessária & Idealmente nula em DC \\
        \hline
        Aplicação digital & Menos comum atualmente & Amplamente utilizado \\
        \hline
    \end{tabular}
\end{table}

A elevada impedância de entrada dos MOSFETs permite reduzir significativamente a potência necessária para controlar seus estados. Além disso, sua estrutura facilita a fabricação de circuitos integrados com milhões ou bilhões de transistores.

Por esse motivo, os MOSFETs são predominantes na construção de processadores, memórias e outros circuitos digitais modernos.

Isso não significa que os BJTs sejam obsoletos: eles continuam sendo amplamente utilizados em amplificadores analógicos, circuitos de potência e aplicações específicas nas quais suas características elétricas são vantajosas.